\section*{Chapter \ref{ch:b3:field-lab-methods-3} --- Field and Lab Methods III}
% Values from the Book 5 scratch calculator calc/ch31.py.

\begin{solution}{exo:b3:field-lab-methods-3:1}
(a)~$\qty{0.01}{\milli\gal} = \qty{e-7}{m.s^{-2}}$, and
$\num{e-7}/\num{9.81} = \num{1.0e-8}$: one part in a hundred million.

(b)~$0.01/0.3086 = \qty{0.032}{m}$, about \qty{3.2}{cm}.

(c)~Together the two effects change gravity by
\qty{0.1967}{\milli\gal.m^{-1}}, the difference $0.3086 - 0.1119$, so
the height must change by $0.01/0.1967 = \qty{5.1}{cm}$.
\end{solution}

\begin{solution}{exo:b3:field-lab-methods-3:2}
$\Delta m = \tfrac43\pi(\qty{50}{m})^3\times\qty{500}{kg.m^{-3}} =
\qty{2.62e8}{kg}$. The peak is $G\Delta m/z^2 =
\num{6.674e-11}\times\num{2.62e8}/\num{e4} = \qty{1.75e-6}{m.s^{-2}}$,
that is \qty{0.175}{\milli\gal}. The half-width is $x_{1/2} = z/1.305 =
\qty{77}{m}$.
\end{solution}

\begin{solution}{exo:b3:field-lab-methods-3:3}
(a)~With $AM = MN = NB = a$, $K = 2\pi/(1/a - 1/2a - 1/2a + 1/a) =
2\pi a = \qty{62.83}{m}$, and $\rho_a = 62.83\times0.239/0.1 =
\qty{150}{\ohm\metre}$. (b)~$K = \pi(15^2 - 5^2)/(2\times5) = \pi\times
200/10 = \qty{62.83}{m}$. The electrodes are then at $0$, $10$, $20$
and \qty{30}{m}, in the order AMNB: it is the Wenner array of spacing
\qty{10}{m}.
\end{solution}

\begin{solution}{exo:b3:field-lab-methods-3:4}
(a)~$V = 0.2998/\sqrt5 = \qty{0.134}{m.ns^{-1}}$. (b)~$t =
2\times10/0.134 = \qty{149}{ns}$. (c)~$\delta = 0.0053\sqrt5/\num{e-3} =
\qty{11.9}{m}$, and a useful depth of about \qty{36}{m}. (d)~In the clay
$\delta = 0.0053\sqrt{20}/0.01 = \qty{2.4}{m}$, a useful depth of about
\qty{7}{m}, and a velocity of only \qty{0.067}{m.ns^{-1}}. The
approximation still holds there: at \qty{100}{MHz},
$\omega\varepsilon = 2\pi\times\num{e8}\times20\times\num{8.85e-12} =
\qty{0.11}{S.m^{-1}}$, ten times $\sigma$.
\end{solution}

\begin{solution}{exo:b3:field-lab-methods-3:5}
(a)~The cylinder has $\lambda = \pi R^2\Delta\rho$ per unit length, and
outside it its attraction is that of a line mass on its axis (Gauss's
theorem on a coaxial cylinder): $\Delta g = 2G\pi R^2\Delta\rho\,z/(x^2
+ z^2)$. Half the peak is reached when $x^2 + z^2 = 2z^2$, so $z =
x_{1/2}$. (b)~The sphere is at $1.305\times100 = \qty{130}{m}$, the
cylinder at \qty{100}{m}: the same profile, read as a sphere, places the
body \qty{30}{\percent} deeper. A sphere's anomaly decays faster with
distance, so a sphere must be deeper to be as wide.
\end{solution}

\begin{solution}{exo:b3:field-lab-methods-3:6}
(a)~Each point stands for $\qty{100}{m}\times\qty{100}{m} =
\qty{e4}{m^2}$, so $\iint\Delta g\,\mathrm{d}A \approx
\num{-438e-5}\times\num{e4} = \qty{-43.8}{m^3.s^{-2}}$, and $\Delta m
= -43.8/(2\pi\times\num{6.674e-11}) = \qty{-1.04e11}{kg}$. (b)~$1.04/1.74
= \qty{60}{\percent}$. The rest is in the flanks of the anomaly outside
the map: the granite reaches a depth of more than a kilometre, and the
anomaly of its deep part decays slowly with distance. Extending the map
far enough, or a model, recovers it.
\end{solution}

\begin{solution}{exo:b3:field-lab-methods-3:7}
(a)~For a thin dyke the anomaly falls to half its peak at $x = \pm z$:
its full width at half the peak is $2z = \qty{60}{m}$, and three
readings across it need a spacing of at most \qty{20}{m}. (b)~$2\times
220/3 = \qty{147}{m}$. (c)~The station spacing needed scales with the
depth of the shallowest body of interest: a survey designed for the deep
dyke, at \qty{150}{m}, would alias the anomaly of the shallow one, which
is only \qty{60}{m} wide; one designed for the shallow body sees the
deep one too, with many more readings than it needs.
\end{solution}

\begin{solution}{exo:b3:field-lab-methods-3:8}
(a)~$F = 80/8 = 10$ and $\phi = F^{-1/m} = 10^{-1/2} = \num{0.316}$.
(b)~$10^{-1/1.8} = \num{0.278}$ and $10^{-1/2.2} = \num{0.351}$. A change
of \num{0.2} in $m$ moves the porosity by about four percentage points;
the resistivities, measured to a few per cent, move it much less. The
cementation exponent controls the uncertainty, and it must be calibrated
on samples of the same rock.
\end{solution}

\begin{solution}{exo:b3:field-lab-methods-3:9}
(a)~At $s = 2h_1$, $(2nh_1/s)^2 = n^2$, so the terms are $k^n(1 +
n^2)^{-3/2}$: $0.3536k$, $0.0894k^2$, $0.0316k^3$. For $k = 0.5$:
$\num{0.1768}$, $\num{0.0224}$, $\num{0.0040}$, and $\rho_a/\rho_1 \approx
1 + 2\times0.2031 = \num{1.406}$ (the full sum gives \num{1.409}). For $k
= -0.5$ the terms alternate in sign, $-0.1768 + 0.0224 - 0.0040$, and
$\rho_a/\rho_1 \approx \num{0.683}$ (full sum \num{0.685}). (b)~For $s \ll
h_1$ every term vanishes and $\rho_a \to \rho_1$; for $s \gg h_1$ every
bracket tends to~1 and $1 + 2\sum k^n = 1 + 2k/(1 - k) = (1 + k)/(1 -
k) = \rho_2/\rho_1$. (c)~$\rho_2/\rho_1 = (1 + k)/(1 - k) = 3$ and $1/3$.
\end{solution}

\begin{solution}{exo:b3:field-lab-methods-3:10}
(a)~In the order AMNB with spacing $a$: $AM = a$, $BM = 2a$, $AN = 2a$,
$BN = a$, so the bracket is $2/a - 1/a = 1/a$ and $K = 2\pi a$.
(b)~With A and B at $\mp s$ and M and N at $\mp l$: $AM = BN = s - l$,
$BM = AN = s + l$, the bracket is $2/(s - l) - 2/(s + l) = 4l/(s^2 -
l^2)$, and $K = \pi(s^2 - l^2)/(2l)$. When $l \ll s$, $\Delta V_{MN}
\approx 2l\,E$ with $E = 2(-\mathrm{d}V/\mathrm{d}r)$ at the centre, the
field of both current electrodes; then $\rho_a = K\Delta V/I \to
\pi s^2/(2l)\times 4l(-\mathrm{d}V/\mathrm{d}r)/I = 2\pi s^2(-
\mathrm{d}V/\mathrm{d}r)/I$. (c)~The ratio of the approximate to the
exact factor is $s^2/(s^2 - l^2) = 1/(1 - 0.01) = 1.010$: one per cent.
\end{solution}

\begin{solution}{exo:b3:field-lab-methods-3:11}
(a)~Reading: \qty{0.010}{\milli\gal}. Height:
$0.1967\times0.05 = \qty{0.010}{\milli\gal}$. Density: the Bouguer
correction is $0.04191\,h\,d$ with $d$ in \unit{g.cm^{-3}}, so
$0.04191\times50\times0.1 = \qty{0.210}{\milli\gal}$. Terrain:
\qty{0.200}{\milli\gal}. Total: $\sqrt{0.010^2 + 0.010^2 + 0.210^2 +
0.200^2} = \qty{0.29}{\milli\gal}$. (b)~The density and the terrain
dominate, each about twenty times the reading error. To detect
\qty{0.1}{\milli\gal}, the density of the slab must be measured on
samples or fitted to the data (the value that least reflects the
topography), and the terrain corrected from a fine elevation model.
(c)~$\sqrt{0.010^2 + 0.010^2 + 0.210^2} = \qty{0.21}{\milli\gal}$.
\end{solution}

\begin{solution}{exo:b3:field-lab-methods-3:12}
(a)~The mass that fits the central reading is $\Delta m = \Delta g(0)z^2/G
= \num{-4.16e-5}z^2/\num{6.674e-11}$: \qty{-5.61e10}{kg} at
\qty{300}{m}, \qty{-9.97e10}{kg} at \qty{400}{m}, \qty{-1.56e11}{kg} at
\qty{500}{m}. The predicted readings at $x = \pm300$ and $\pm600$ are
\num{-1.47} and \num{-0.37}, \num{-2.13} and \num{-0.71}, \num{-2.62}
and \qty{-1.09}{\milli\gal}. The \omterm{def:b3:field-lab-methods-3:inversion-practice}{misfits} are $\chi^2 = 1238$, $3.7$ and
$856$. (b)~Only the model at \qty{400}{m}, with $\chi^2 = 3.7$ for five
data, fits within the errors. (c)~Yes: any spherically symmetric body
with the same mass and centre gives exactly the same anomaly, and other
shapes can come within the errors. A density measured on samples, a
borehole or another method would decide.
\end{solution}

\begin{solution}{pb:b3:field-lab-methods-3:1}
\textbf{1.} $(2846.708 - 2846.512)/\qty{2.0}{h} = \qty{0.098}{\milli\gal.h^{-1}}$.

\textbf{2.} $2851.075 - 0.098\times0.5 - 2846.512 = \qty{4.514}{\milli\gal}$;
$2850.497 - 0.088 - 2846.512 = \qty{3.897}{\milli\gal}$;
$2849.392 - 0.127 - 2846.512 = \qty{2.753}{\milli\gal}$.

\textbf{3.} The heights above the base are $-27.1$, $-30.6$ and
\qty{-18.1}{m}; adding $0.1967\,\Delta h$ gives the Bouguer anomalies
$\num{-0.817}$, $\num{-2.122}$ and \qty{-0.808}{\milli\gal}. The
stations are lower than the base and nearer the centre of the earth,
so their raw readings are higher; the correction removes that and
leaves a low over the granite.

\textbf{4.} The extra $0.098 - 0.040 = \qty{0.058}{\milli\gal.h^{-1}}$,
\qty{0.12}{\milli\gal} over the loop, is the solid-Earth tide. Over two
hours it changes almost linearly, so a straight line through the two
base readings removes it with the drift.

\textbf{5.} The errors are $\num{0.008}$, $\num{0.009}$ and
\qty{0.017}{\milli\gal}. Rounding each height to \qty{0.1}{m} can make
a height difference wrong by up to \qty{0.1}{m}, that is
\qty{0.020}{\milli\gal}: the rounding explains the errors, the tide
left by the straight line being below \qty{0.001}{\milli\gal}.

\textbf{6.} $0.01/0.1967 = \qty{5.1}{cm}$. The \qty{3.3}{cm} of the manuals is
$0.01/0.3086$: the free-air anomaly does not depend on the density
chosen for the Bouguer correction, and the height must serve it too.

\textbf{7.} $x_{1/2} = (3.952 - 3.048)/2 = \qty{0.452}{km}$, and $z = 1.305
\times452 = \qty{590}{m}$ below the ground.

\textbf{8.} $\Delta m = \Delta g\,z^2/G = \num{-2.293e-5}\times590^2/
\num{6.674e-11} = \qty{-1.20e11}{kg}$. With $\Delta\rho =
\qty{-130}{kg.m^{-3}}$ its volume is \qty{9.2e8}{m^3} and its radius
\qty{603}{m}, larger than the depth of its centre: the sphere would
stick out of the ground. No sphere with the granite's density contrast
fits; the body must reach close to the surface.

\textbf{9.} The granite runs from the surface to a great depth. Its shallow
part, near the stations, makes the peak narrow, and the rule, which
knows only the width, puts a single sphere at a moderate depth; the
deep part adds broad flanks that the sphere lacks (the figure of the
profile), and the mass they hold is missed.

\textbf{10.} $2\pi G\Delta\rho = 2\pi\times\num{6.674e-11}\times(-130) =
\qty{-5.45e-8}{s^{-2}}$. With $z_1 = 0$, $z_2 = \qty{1645}{m}$, $a =
\qty{461}{m}$: the bracket is $1645 + 461 - 1708.4 = \qty{397.6}{m}$,
and $\Delta g = \qty{-2.17}{\milli\gal}$. For an unlimited depth the
bracket tends to $a$, and $\Delta g = \qty{-2.51}{\milli\gal}$.

\textbf{11.} The bracket must be $2.29/\num{5.45e-3} = \qty{420}{m}$, which a
floor at about \qty{2.6}{km} gives. Any floor between about
\qty{1.6}{km} and an unlimited depth changes the peak by a few tenths
of a milligal only, and a slightly wider or denser body gives the same
peak with a shallower floor: the floor is not resolved.

\textbf{12.} With $z_1 = \qty{6}{m}$: $1639 + 461.0 - 1708.4 = 391.6$, and
$\Delta g = \qty{-2.14}{\milli\gal}$, \qty{0.03}{\milli\gal} less than
with the top at the ground. That is below the terrain and density
errors of the survey: gravity cannot tell whether the granite reaches
the surface under the alluvium. The sounding can.

\textbf{13.} $K = \pi(20^2 - 1^2)/(2\times1) = \qty{626.7}{m}$, and $\Delta V =
\rho_a I/K = 821\times0.050/626.7 = \qty{65.5}{mV}$.

\textbf{14.} $V = 2\times1/\sqrt{17.1^2 - 10.7^2} = 2/13.34 =
\qty{0.150}{m.ns^{-1}}$, $\varepsilon_r = (0.2998/0.150)^2 = 4.0$, a
dry sand; the pipe is $0.150\times10.7/2 = \qty{0.80}{m}$ deep.

\textbf{15.} $0.150\times20.0/2 = \qty{1.50}{m}$.

\textbf{16.} $20/821 = \num{0.0244}$, $30/1309 = \num{0.0229}$, $40/1684 =
\num{0.0238}$: $S \approx \qty{0.0237}{S}$.

\textbf{17.} The dry layer contributes $1.5/600 = \qty{0.0025}{S}$, the
saturated one $\qty{0.0212}{S}$, so its thickness is $0.0212\times222 =
\qty{4.7}{m}$, and the granite lies $1.5 + 4.7 = \qty{6.2}{m}$ below
the floor. The estimate is high: the granite is not infinitely
resistive, so at these spacings $\rho_a$ still lies below the line
$s/S$, and $s/\rho_a$ overestimates $S$. (The region's model has
\qty{6}{m} of alluvium.)

\textbf{18.} $\phi = (20/222)^{1/2} = \num{0.30}$. With $m = 1.3$, $\phi =
0.0901^{1/1.3} = \num{0.16}$; with $m = 2.8$, $\phi = \num{0.42}$.

\textbf{19.} Subtracting the base from the line gives, from west to east:
3.4, 4.1, 4.8, 5.9, 7.4, 9.2, 11.5, 14.2, 17.2, 19.8, 21.0, 20.1, 17.6,
14.5, 11.5, 9.1, 7.2, 5.7, 4.6, 3.8 and \qty{3.1}{nT}.

\textbf{20.} Uncorrected, the peak is \qty{28.0}{nT} at \qty{2.20}{km}, and
half of it, \qty{14.0}{nT}, is never reached to the east: the line ends
at \qty{15.2}{nT}, because the field rose by \qty{12}{nT} during the
walk.

\textbf{21.} The peak is \qty{21.0}{nT} at \qty{2.20}{km}; half of it is
\qty{10.5}{nT}. West: between \num{9.2} at \num{1.95} and \num{11.5} at
\qty{2.00}{km}, at $1.95 + 0.05\times1.3/2.3 = \qty{1.978}{km}$. East:
between \num{11.5} at \num{2.40} and \num{9.1} at \qty{2.45}{km}, at
$2.40 + 0.05\times1.0/2.4 = \qty{2.421}{km}$. The half-width is $(2.421
- 1.978)/2 = \qty{0.221}{km}$.

\textbf{22.} A symmetric anomaly in a field inclined at \ang{65} means that only
the vertical magnetisation acts: the body strikes along the magnetic
meridian, north--south. The rule of
\cref{prop:b3:field-lab-methods-3:dyke} assumes a thin, vertical,
two-dimensional dyke reaching far down, magnetised by induction, whose
top acts as a line of magnetic charge.

\textbf{23.} $\kappa t = 2\pi z\,\Delta T_\mathrm{max}/(F\sin^2 I) = 2\pi\times
221\times21.0/(48000\times0.8214) = \qty{0.74}{m}$. With $\kappa =
0.03$, $t = \qty{25}{m}$. Only the product is fixed: a dyke twice as
thick with half the susceptibility gives the same anomaly.

\textbf{24.} The anomaly is \qty{443}{m} wide at half its peak and the stations
are \qty{50}{m} apart: nine readings across it, three times the rule's
minimum.

\textbf{25.} The top of the dyke is
\textbf{\qty[text-series-to-math=true]{220}{m}} below the ground. The
geological model puts the unconformity \qty{222}{m} down at this point:
the dyke rises to the base of the cover and is cut off there, older
than the Triassic sandstone that buries it.
\end{solution}
